\begin{explanationbox}
把椭圆想成一条闭合的轨道，点 $M$ 固定，
点 $P(\theta)$ 沿椭圆运动一周。
此时 $MP$ 的长度随 $\theta$ 改变，
必定存在最短和最长距离。

\textbf{为什么优先计算距离平方？}
\begin{itemize}
  \item 距离本身含根号，求导不方便；
  \item $MP\ge0$，平方不会改变距离的大小顺序；
  \item 因此先求 $F(\theta)=MP^2$ 的最值，再开方即可。
\end{itemize}

\textbf{直观抓手：}三角参数
$P=(a\cos\theta,b\sin\theta)$ 把“椭圆上任意动点”
化成了“一个参数 $\theta$ 在变化”。
不是对 $x,y$ 分别求最值，也不是直接套用圆的结论。

若 $x_0=0$ 或 $y_0=0$，常可利用
$\sin^2\theta+\cos^2\theta=1$
化成二次函数；
一般位置的外点则通过求导和正切半角代换求解。
\end{explanationbox}
